\section{Introduction}
Economic counterfactuals require more than a numerical solution at a single calibration. A researcher changes an operating benefit, an adjustment cost, or the distribution of a financial shock and asks whether a policy response reflects an economic mechanism rather than approximation error. Repeating an optimization at every parameter value is one response. Reusing an approximate policy is another, but it is useful only when its welfare loss and the scope of its validity can be evaluated. This paper studies that problem for dynamic economies in which control changes both financial resources and the agent's preferences.

An operating mandate also requires an economic account of its obligations. We introduce voluntary participation, an externally funded signing grant, finite enforcement capacity, and a right to surrender after a minimum term. Preference adjustment can lower both the compensation required to operate and the enforcement capacity needed to implement continued operation. Its relative value across initial position classes can nevertheless reverse when the surrender margin changes. We derive that distinction and map the two investment-permission boundaries using an exact first-date breakpoint reduction. These results tie a certified risk comparison to the price of the contract that makes it relevant.

Neural Bellman Operators (NBO) separate policy proposal, evaluation, improvement, and certification. A critic represents continuation value, an actor proposes a feasible action, an improvement calculation checks that action, and an independent evaluation measures the induced policy. We use this separation to develop a counterfactual method with an explicit distinction between changes in rewards and changes in transition laws. The certificate, rather than a training objective or a preferred approximation architecture, determines which economic conclusions are warranted.

The first distinction is mathematical. For fixed transitions, the value of a policy is affine in reward coefficients. Policy caches and optimistic linear support exploit precisely this structure. Once a parameter changes the transition law at every date, even a fixed policy can have a nonconvex value in that parameter. We identify the cross-operator term that makes an uncorrected endpoint chord fail and construct a positive backward correction yielding an upper value throughout a parameter interval. Its additional error is second order in interval width under explicit Lipschitz conditions. A matched comparison shows that a localized count-information upper value is sharper, whereas the correction compresses triangular horizon dependence into a linear number of action-wide kernel applications. We give a total coefficient-error bound and a finite-refinement guarantee that remain valid when optimal actions switch. Independently evaluated policy values have a finite Bernstein representation for an affine mixture of transition laws. Combining their maximum with the corrected upper value and statewise policy switching gives a welfare certificate at every state and date over a continuum of transition counterfactuals. The result also admits approximate neural endpoint values after a certified residual correction.

The second distinction is economic. In a portfolio economy with adjustable risk aversion and first-exit liquidation, a constant operating benefit $m$ and a liquidation increment $\ell$ affect behavior through the contrast $d=m-\rho\ell$. The annuity identity explains which contract changes matter, but not whether a portfolio reversal is caused by preference formation. We partition policies by the sign of the risky position held over a positive-length initial rebalancing interval. The slope of a regular indifference locus with respect to adjustment cost equals the difference in adjustment effort between the two risk classes divided by their difference in discounted duration. A separate option-value decomposition compares deliberate preference adjustment with the otherwise identical economy in which adjustment is unavailable. It identifies the quantity that must change the ranking, but an identity alone does not predict that change. We derive a primitive sufficient condition from consumption exposure and the curvature of CRRA utility with respect to its preference index. It predicts a strictly positive risk-specific adjustment premium over an entire probability--cost family. In the full finite stopped economy, class-specific tensor certificates then establish opposite risk rankings with and without adjustment over a joint region of contract and transition parameters.

The finite implementation retains consumption, risky investment, costly preference adjustment, stochastic preference and wealth shocks, their covariance, and the original settlement rule. It includes the union of both historical action meshes and the frozen neural proposals. The reward-transfer bank retains its all-state, all-date $10^{-3}$ welfare certificate and its separate first-action sign exclusions. The new uncertainty experiment changes the probability of two one-step shock laws, each with its own killed transition kernel, and certifies their entire mixture interval. This changes the dynamics; it is not a relabeling of a reward coefficient. The finite-chain mixture is distinguished from changing a Brownian correlation parameter without a discretization bound.

The economic interpretation is developed further by specifying the decision calendar and settlement lottery, varying its wealth granularity with both adjustment regimes recomputed, and separating the resulting movement of a risk boundary from the relative adjustment option. We retain the original region, including the loss of its no-adjustment sign at a finer-settlement corner. A common nested region is additionally certified for each of three finite protocols. A preference over preference states supplies the cardinal primitive not identified by a CRRA index alone. We prove a uniform option bound under bounded common cardinal tilts and then use full dynamic reoptimization to distinguish that primitive from operating-duration, liquidation, and covariance incentives. These results concern explicit economic experiments rather than a sign assigned to an unspecified continuous-time limit.

The computational comparisons assign credit to the object doing the work. We execute an exact-oracle scalar optimistic-linear-support comparison, distinguish an initial-distribution objective from simultaneous statewise coverage, and remove all neural actions from the lower policy bank while retaining the original full-menu upper benchmark. In the resource economy, we add an analytically assembled, probability-preconditioned parametric quadratic program that reuses the same structure across fresh queries. It is substantially faster than the learned arm in this short quadratic model. This comparison changes the inference from the earlier crossover: nonlinear continuation can outperform the fitted quadratic approximation without outperforming the exact reusable quadratic program. The transition-law experiments compare upper oracles on identical anchors, feasible policies, and subdivision cells, charging shared setup and evaluation work to each arm. They also measure autonomous dyadic refinement from empty policy caches. The compressed and count constructions occupy different points on a precision--work frontier; the rectangular relaxation is a further regional comparator. The economic use is a decision-specific region certificate rather than a generic welfare bound placed beside a separately optimized risk bracket. A new streaming comparison charges both upper methods for the same model construction, full-menu endpoint optimization, feasible-policy evaluation, and independent replay at that decision target. It also tests a mesh-only feasible lower bank against the unchanged full target. This comparison attributes a gain to compression without treating learned proposals as indispensable.

The general error analysis and the broader economic extensions remain part of the argument. A finite-horizon theorem distinguishes evaluation, improvement, boundary, and transition errors. Exact policy improvement provides a conditional benchmark; counterexamples exclude identification of attracting actor stationary points with globally optimal policies. Merton and Epstein--Zin examples test actual evaluation and control updates. Sophisticated temporal selves require a different improvement criterion from ordinary continuation evaluation. Persistent capacity competition tests unilateral dynamic deviations at every date and state. These exercises retain their separate economic purposes rather than becoming an aggregate count of successes.

\subsection{Related Literature and the Contribution}
Policy improvement and approximate dynamic programming precede neural approximation. \citet{jacka2015} study continuous-time improvement under explicit control hypotheses. The nearest neural approaches also combine continuation approximation and policy improvement: fixed-policy PDE evaluation in \citet{kim2025,kim2026}, operator learning in \citet{lee2024}, martingale evaluation in \citet{cai2025}, and direct control approximation in \citet{han2016}. Neural evaluation, separate parameter blocks, and avoidance of a closed-form maximizer are not claims of priority here. Our operator assumptions and error accounts specify when such computations support the economic use of a policy.

Successor features and generalized policy improvement \citep{barreto2017}, policy switching \citep{chang2021}, and policy-cache upper bounds \citep{nemecek2021} provide the reward-transfer foundations. \citet{alegre2022} combine optimistic linear support with successor features and generalized policy improvement. Their work is a necessary comparator: selecting an unsatisfactory reward weight, solving there, and adding its policy is not a new algorithm merely because the reward weight is a contract. We prove below that, with exact statewise endpoint optima and one policy per scalar anchor, our adjacent-policy envelope equals the full scalar-OLS policy envelope. We measure the cost of strengthening an initial-distribution guarantee to simultaneous statewise coverage, but do not present that strengthening alone as a new transfer principle. Our switching rule also differs from full action-value generalized policy improvement; the comparison does not identify the two operators.

The additional transfer result concerns \emph{changing transition laws}, outside the fixed-feature reward correspondence. Lipschitz dependence of value on rewards and transitions is established in the changing-environment literature \citep{csaji2008}. Theorem \ref{thm:corrected_chord} instead gives an explicit endpoint-supersolution correction involving the product of a kernel difference and a continuation-value difference. Its implementation is paired with exact feasible-policy polynomials, a statewise welfare bound, and an independently computed information-relaxation upper value. The latter is an instance of the general approach of \citet{brown2010}, not a new duality theorem. Regional verification and parameter lifting \citep{quatmann2016,heck2025} are also direct comparators for repeated transition parameters; we distinguish their dependence relaxations and implement a finite-horizon rectangular upper recursion on the same economic target. Section \ref{sec:r8_correspondence} gives the date-augmented reward-MDP correspondence, the order of action and endpoint information, the treatment of killed mass and terminal rewards, and the precise role of signed compression. The economic switching-locus result uses the envelope logic of \citet{milgrom2002}; the envelope identity organizes the calculation but does not itself establish a duration ordering. The additional economic result is a primitive consumption-exposure condition and an integrated class-value certificate, not a claim to a new envelope theorem.

The broader computational literature includes recursive and projection methods \citep{stokey1989,judd1998}, sparse grids \citep{brumm2017}, and neural differential-equation solvers \citep{raissi2019,han2018}. Our input-convex critic follows \citet{amos2017}. Riccati recursion and a reusable convex quadratic program are the appropriate structural comparators for the corresponding models. We retain both, including results unfavorable to learned deployment. Thus the paper's incremental case rests on the controlled compression of transition-law upper information, the resulting precision--work and anchor-demand comparisons, and a certified preference-adjustment effect on economic choice. The architecture proposes policies; theorems and independent class evaluation determine which counterfactual conclusions those policies support.
